% Transcription — fonds Alexandre Grothendieck, shelfmark 32, batch 2 (pages 21-40).
% Established from the facsimile archives/batches/32/batch-02.pdf (a mirror of
% the folder's PDF; no cover sheet, so PDF page k = folder page 20+k), per the
% skill transcribe-grothendieck.
%
% Pass: Opus 5.5 (claude-opus-5-5), 2026-09-24 — first pass, unchecked against the pages by a human.
%
% Correction material for SGA 4, exposés XVII and XVIII, in several hands.
% The file does not say whose hand is whose, except where a page is signed
% (page 30, « P. Deligne »). Pencil annotations by a reader on inked text go
% into \marginal{} or \add{}, as in batch 3.
%
%   Pages 21-23 — three more correction slips like pages 15-20 of batch 1
%     (« XVII -- 1.2 (4) », « XVIII.1.2 (7 », « XVIII 1.2 6 »): lines to
%     change, « remplacer 1.2.2.1 par (1.2.3) et (1.2.4.2) », a struck
%     « Lie Pic_X = H^1(X, O_X) ? ».
%   Page 24 — ink, « XVIII -- 1.2 »: for p : P(E) → S, the equivalence of
%     I (G ≃ p_* p^* G and Hom(G_m, G) ≃ R^1 p_* G) and II (the functor
%     (1.2.1.1) from G-torsors on S to those on P(E) is an equivalence),
%     with three remarks (the exact sequence 0 → H^1(S,G) → H^1(P(E),G) →
%     Γ(S, R^1 p_* G) → 0 given a section, locality on the base,
%     injectivity) and the proofs of both implications.
%   Pages 25-27 — three leaves of a redaction of XVIII 1.2 (blue ink, squared
%     paper, folios 8 and 12 on 26-27): Lemme 1.2.9 (S local artinian,
%     f : X → S proper and smooth, G smooth commutative; if H^0(X_s̄, O) =
%     k(s̄), Pic^τ X_s̄ = 0 and G_s̄ contains no G_m, then f^* is an
%     equivalence of torsor categories) and its proof (a: locality and full
%     faithfulness; b: the abelian-variety part via Raynaud and torsion; c:
%     deformation from S_0 along an ideal of square zero, classified by
%     H^1(X, f^* Hom(L, I)) = 0); page 27, struck through by two diagonals,
%     proves (1) ⇒ (1') for finite locally free G via Cartier duality and
%     the case (iii) via 0 → G_m^n → G_0 → H → 0.
%   Page 28 — blue ink, upside down, struck by diagonals: brief notes on
%     complex structures on the disc, flat PSL_2(R)-torsors on a curve,
%     whether they lift to SL_2(R), the obstruction in H^2(X, Z/2), and the
%     relation among lifted generators (x_1 x_2)...(x_{2g-1} x_{2g}) = ±1.
%   Page 29 — pencil: references to check in SGA 4 XVII (XVII 0.13, sheaves
%     prime to the residue characteristics), « 1.2.9 : plat / lisse par
%     Illusie ».
%   Pages 30-31 — a letter signed P. Deligne (« L 4 »): what is missing
%     (Gysin, g^* Rf^! ≃ Rf'^! g^*, compatibilities, local fundamental
%     classes), what he will and will not fill before « SGA 4 ne devienne un
%     livre », the list of enclosures; then the « Liste de compatibilités
%     manquantes » a)-i) the letter announces.
%   Pages 32-33 — two slips headed in red « ETAT Avant Illusie »: the
%     low-degree Leray sequence, the pairing ⟨O(E), O(F)⟩ = det Rf_*(O_E ⊗
%     O_F)^{-1} and norms; divisors and the tame symbol
%     (-1)^{v f v g} g^{v f}/f^{v g}; Sym of a presentation T'' ⇉ T' → T and
%     descent (« desc. el. », « desc. finie »).
%   Page 34 — red: « ETAT PRIMITIF. contient : ? compatibilité … ».
%   Page 35 — blue ink, struck by diagonals: sketches of the adjunctions
%     f_! ⊣ f^!, the two base-change squares (« dém : fib + cofib »), and
%     Hom(f_! K, L) ~ Hom(K, f^! L) with a variant.
%   Pages 36-40 — blue ink fair copy of part of a § 3.1 on Rf^!, annotated
%     in pencil: end of 3.1.6 (the base-change isomorphism (3.1.6.3)
%     Rg'_* Rf'^! ≃ Rf^! Rg_* and the four base-change arrows (3.1.6.4));
%     page 37 a struck heading « 3.1.3 Constructions du foncteur Rf^! »;
%     3.1.7 (the projection-type arrow (3.1.7.1), its transpose (3.1.7.2),
%     the induction homomorphism (3.1.7.3) and the diagram proving it an
%     isomorphism), Proposition 3.1.8 (« formule d'induction »), 3.1.9
%     (restriction of scalars commutes with Rf^!: the ring sheaf plays only
%     a dummy role); page 40 a struck « Prenons Soit B ». Batch 3 continues
%     the same § 3.1 (3.1.14 onwards).
%
% Nothing is skipped: every page carries mathematics or correction matter.
% Pages 37 and 40 are given for their single struck line.
%
% The dating is the inventory's own, for the whole shelfmark.
\documentclass[11pt,a4paper]{article}
% Le préambule commun à toutes les transcriptions du fonds Grothendieck.
%
% Il tient en une idée : **l'incertitude se compose, elle ne se résout pas.**
% Une transcription qui lisse un mot douteux a détruit la seule chose pour
% laquelle on la faisait — un lecteur doit pouvoir savoir, à chaque endroit,
% ce qui était sur le feuillet et ce qui a été ajouté. Les six macros qui
% suivent sont donc l'appareil critique, et non de la décoration.
%
% Les mêmes commandes sont reconnues par `scripts/render.mjs`, qui produit la
% vue de lecture du site. Une macro ajoutée ici doit l'être aussi là-bas,
% faute de quoi le rendu échouera bruyamment — ce qui est voulu : un rendu
% silencieusement faux est pire qu'un rendu absent.

\ProvidesPackage{grothendieck}[2026/08/08 Transcriptions du fonds Grothendieck]

\usepackage{amsmath,amssymb,amsthm}
\usepackage{mathrsfs}
\usepackage{stmaryrd}
% Les diagrammes commutatifs sont partout dans ce fonds : c'est la langue
% même de la géométrie algébrique de Grothendieck, et une transcription qui
% les rendrait en prose ne transcrirait rien.
\usepackage{tikz-cd}
% Français par défaut — c'est la langue des feuillets — mais l'anglais est
% chargé aussi : la traduction et le résumé sont des documents anglais, et la
% typographie française (espaces avant les deux-points) y serait une faute.
% Ces documents basculent par \selectlanguage{english} en tête de corps.
\usepackage[english,french]{babel}
\usepackage{fontspec}
\usepackage{geometry}
\geometry{a4paper,margin=25mm,marginparwidth=28mm}
\usepackage{marginnote}
\usepackage{xcolor}
% ulem, not soul. `soul`'s \st and \ul are built on a character-by-character
% scan that XeLaTeX's font machinery breaks: the document fails with an
% undefined control sequence rather than degrading. ulem does the same two jobs
% and is Unicode-engine safe. `normalem` stops it hijacking \emph, which the
% transcriptions use for Grothendieck's own emphasis.
\usepackage[normalem]{ulem}
\usepackage{hyperref}
\hypersetup{colorlinks=true,linkcolor=black,urlcolor=black,pdfborder={0 0 0}}

% --- Métadonnées du lot -----------------------------------------------------
% Reprises telles quelles par le script de rendu, qui les affiche en tête de
% la vue de lecture. Elles fixent ce que le fichier prétend couvrir : sans
% elles, une transcription flotte sans point d'attache dans un fonds de seize
% mille feuillets.

\newcommand{\folder}[1]{\gdef\grothFolder{#1}}
\newcommand{\batch}[1]{\gdef\grothBatch{#1}}
\newcommand{\leaves}[2]{\gdef\grothFirst{#1}\gdef\grothLast{#2}}
\newcommand{\foldertitle}[1]{\gdef\grothFoldertitle{#1}}
\gdef\grothFolder{?}\gdef\grothBatch{?}\gdef\grothFirst{?}\gdef\grothLast{?}\gdef\grothFoldertitle{}

% --- Le résumé --------------------------------------------------------------
%
% Il ouvre la lecture modernisée, sous le titre « Résumé », et il est composé
% en petit. Ce n'est pas un ornement : c'est le seul passage du document qui
% ne soit pas des mathématiques, et un lecteur doit pouvoir le reconnaître
% avant de l'avoir lu — puis le sauter s'il connaît déjà le sujet, ou s'y
% arrêter s'il ne le connaît pas. Le filet qui le suit marque la reprise du
% corps.

\newenvironment{resume}
  {\small\setlength{\parskip}{0.45em}}
  {\par\medskip\hrule\medskip}

% --- Mots-clés ---------------------------------------------------------------
%
% En anglais, en clôture du résumé de la lecture modernisée : le vocabulaire
% moderne sous lequel chercher ce que les feuillets construisent. Le manifeste
% du site (`npm run manifest`) les extrait pour étiqueter la cote entière —
% cette ligne est la source unique des tags, il n'y a pas de fichier à côté.
\newcommand{\keywords}[1]{\par\smallskip\noindent
  {\footnotesize\textsc{Keywords} --- \textit{#1}}\par}

% --- L'appareil critique ----------------------------------------------------

% Le numéro de feuillet, en marge. C'est l'ancre qui permet de régler un
% désaccord sur une formule en regardant *un* feuillet plutôt qu'une chemise
% de six cents. Le site s'en sert aussi pour tourner les pages du fac-similé
% quand on fait défiler la transcription.
\newcommand{\leaf}[1]{%
  \marginnote{\footnotesize\color{gray!70}\textsf{#1}}%
  \label{leaf:#1}\ignorespaces}

% Illisible : on ne devine pas. Un mot inventé qui se lit comme les autres est
% le pire résultat possible.
\newcommand{\ill}{\textcolor{red!60!black}{[\,\ldots\,]}}

% Lecture douteuse : le mot est donné, et signalé comme incertain.
\newcommand{\uncertain}[1]{\uline{#1}}

% Ajout de l'éditeur — une lettre manquante, un mot sous-entendu.
\newcommand{\add}[1]{\textcolor{blue!50!black}{[#1]}}

% Biffé par Grothendieck lui-même. Ce qu'il a rayé fait partie du document :
% une rature dit souvent où la pensée a changé de direction.
\newcommand{\struck}[1]{\sout{#1}}

% Note du transcripteur, sur l'état matériel du feuillet.
\newcommand{\note}[1]{\par\smallskip{\small\color{gray!60!black}\textit{[#1]}}\par\smallskip}

% Note marginale de Grothendieck, distincte du corps du texte.
\newcommand{\marginal}[1]{\par\smallskip\noindent\hspace{1em}%
  {\small\color{gray!60!black}#1}\par\smallskip}

% --- Titre ------------------------------------------------------------------

\AtBeginDocument{%
  \begin{center}
    {\large\bfseries Fonds Alexandre Grothendieck --- cote n\textsuperscript{o}~\grothFolder}\\[2pt]
    {\itshape\grothFoldertitle}\\[4pt]
    {\small feuillets \grothFirst--\grothLast\ (lot \grothBatch)}\\[2pt]
    {\footnotesize\color{gray!60!black}Transcription établie d'après le fac-similé numérisé
     par l'Université de Montpellier.\\
     Les lectures incertaines sont soulignées, les illisibles notées [\,\ldots\,],
     les ajouts entre crochets.}
  \end{center}
  \vspace{1em}}

\endinput


\folder{32}
\batch{2}
\pages{21}{40}
\dating{[vers 1963-1977]}
\watermark{Édition de démonstration}
\foldertitle{Corrections (références, rajouts) [dont SGA 4] : tapuscrit (s.d.), notes manuscrites (s.d.)}

\begin{document}

\section*{Fiches de corrections à XVIII 1.2, suite (pages 21 à 23)}

\page{21}
\note{Fiche au crayon, du même type que celles des pages 15 à 20 (lot 1). En
tête « XVII -- 1.2 », avec un « 4 » cerclé à droite ; la dernière ligne, à
l'encre bleue, dit XVIII.}

XVII -- 1.2 \quad (4)

($\ell$.\ 9) du haut

cf.\ (1.2.2.4)

mettre (1.2.4.1) \note{« ok » dans la marge, d'une autre écriture.}

\emph{stop p.\ 4.} XVIII.1-2. \note{Cette ligne est à l'encre bleue, sous un
trait bleu.}

\page{22}
\uncertain{XVIII}.1.2 \quad (7

\uncertain{à préciser} $R^1 f_{*}(\mathbb{G}_m)$
\note{les deux premières lignes sont à l'encre bleue ; le reste de la fiche
est au crayon, sous un trait.}

ligne \uncertain{0} du bas \note{le chiffre est un petit cercle : 0 ou 6.}

remplacer 1.2.2.1 par (1.2.3) et (1.2.4.2)

\page{23}
XVIII.\ 1.2 \quad 6

\struck{p.\ 1.2.3}

$\ell$.\ 4 du bas \note{« ok » dans la marge.}

remplacer 1.2.2.1 par (1.2.3)

\struck{$\operatorname{Lie} \operatorname{Pic}_X = H^1(X, \mathcal{O}_X)$ ?}
\note{cette ligne, isolée entre deux traits, est barrée d'un trait ondulé.}

\section*{XVIII 1.2 : torseurs sur $P(E)$ (page 24)}

\page{24}
\note{Feuillet à l'encre noire, sur papier uni ; en tête « XVIII -- 1.2 »,
souligné deux fois. $p : P(E) \to S$ est le fibré projectif, et (1.2.1.1)
le foncteur des $G$-torseurs sur $S$ vers ceux sur $P(E)$ dont il est
question dans les feuillets qui suivent.}

\begin{itemize}
\item[I] I.a\quad $G \to p_{*} p^{*} G$

I.b\quad $\underline{\operatorname{Hom}}(\mathbb{G}_m, G) \xrightarrow{\ \sim\ } R^1 p_{*} G$
\item[II] le foncteur (1.2.1.1) est une équivalence.
\end{itemize}
\note{la flèche de I.a ne porte pas de tilde sur la page ; celle de I.b en
porte un.}

\emph{Remarques 1} : \struck{disparaît} moyennant (I.a) et l'hypothèse d'une
section on a une suite exacte
\[
0 \to H^1(S, G) \to H^1(P(E), G) \xrightarrow{\ \theta\ } \Gamma(S, R^1 p_{*} G) \to 0
\]

\emph{Remarque 2} : La pleine fidélité de (1.2.1.1) est locale sur la base et
moyennant la pleine fidélité, la surjectivité essentielle est locale sur la
base.

\emph{Remarque 3} : I.b (et I.a) entraîne que $\underline{\operatorname{Hom}}(\mathbb{G}_m, G)
\to H^1(P(E), G)$ est injectif.

\struck{I) $\Rightarrow$ II)} \struck{pleine fidélité} : \uncertain{on se
ramène au cas où} il y a une section.

I $\Rightarrow$ II.\quad \emph{pleine fidélité} : si $K(\kappa_0, \varphi) =
K(\kappa'_0, \varphi')$ alors $\varphi = \varphi'$. Alors $\kappa_0 =
\kappa'_0$ et I.a entraîne la pleine fidélité.

\emph{essentielle surjectivité} : Soit $K \in H^1(P(E), G)$ ; $\theta(K) =
\varphi(\mathcal{O}(1))$ pour un certain $\varphi$ et $\theta(K) -
\varphi(\mathcal{O}(1)) \in H^1(S, G)$.
\note{tel qu'écrit : c'est sans doute la différence de $K$ et du torseur
défini par $\varphi(\mathcal{O}(1))$ qui est visée, $\theta$ étant à valeurs
dans $\Gamma(S, R^1 p_{*} G)$.}

II) $\Rightarrow$ I)\quad I-a) résulte de la pleine fidélité.

I.b) \struck{$\underline{\operatorname{Hom}}(\mathbb{G}_m, G) \to$} trivial
\uncertain{vu ce qui précède}.

\section*{XVIII 1.2 : feuillets d'une rédaction (pages 25 à 27)}

\page{25}
\note{Encre bleue, sur papier quadrillé, d'une écriture soignée ; les
insertions sont à la même encre, une flèche et une accolade à l'encre verte.
Feuillet sans numéro ; les pages 26 et 27 portent « XVIII 1.2 » et les
folios 8 et 12.}

\textbf{Lemme 1.2.\struck{\uncertain{7}}9.} Soient $S$ un schéma local
artinien, \add{$\bar s$ un point géométrique de $S$}, $f : X \to S$ un
morphisme propre \add{et lisse} \struck{\ill{}}, \struck{\ill{}} et $G$ un
schéma en groupes commutatif lisse de type fini sur $S$. Si $H^0(X_{\bar s},
\mathcal{O}) = k(\bar s)$ et $\operatorname{Pic}^{\tau} X_{\bar s} = 0$,
\add{et} \struck{alors le foncteur} si $G_{\bar s}$ ne contient pas de
sous-groupe $\mathbb{G}_m$, alors le \struck{le} foncteur $f^{*}$ de la
catégorie des $G$-torseurs sur $S$ dans celle des $f^{*}G$-torseurs sur $X$
est une équivalence.
\note{le numéro du lemme est corrigé : un chiffre biffé, lu 7, puis 9. Le
passage biffé après « propre » (deux ou trois mots) est illisible sous les
ratures ; « $\bar s$ un point géométrique de $S$ » et « et lisse » sont
ajoutés au-dessus de la ligne.}

b.\ Si $S$ est le spectre d'un corps \struck{séparablement clos} et si $A$
\add{est} une variété abélienne, alors $H^1(X, A)$ est de torsion (Raynaud
[1, XIII 2.6]), donc réunion des images des groupes $H^1(X, {}_nA)$
\struck{; ceux-ci sont nuls par 1.2.2.7}. Si $G$ est extension d'une variété
abélienne $A$ par un groupe \struck{\ill{}} affine $G_0$, \struck{$H^1(X,G)$
est donc}
\note{une flèche verte part de l'alinéa a.\ ci-dessous et pointe vers
« b. » : l'ordre des deux alinéas est sans doute à inverser.}

a.\ Le problème posé est local sur $S$ pour la topologie étale ; on se ramène
donc au cas $S$ strictement hensélien, donc au cas où \struck{l'on peut donc
supposer que} $X/S$ admet une section ; le foncteur $f^{*}$ est alors
pleinement fidèle d'après \struck{1.2.2.4}. 1.2.3 et (1.2.4.2)
\note{l'alinéa a.\ est tenu, dans la marge, par une accolade verte.}

\page{26}
\note{Folio 8 et « XVIII 1.2 » à l'encre rouge. Suite de l'alinéa b.\ de la
page précédente.}

et si $G^n$ est l'image réciproque de ${}_nA$, le groupe $H^1(X, G)$ est
réunion des groupes $H^1(X, G^n)$. Les groupes $G^n$ sont affines (non
nécessairement lisse) et, sous l'hypothèse faite, l'assertion résulte de
1.2.\struck{7}8. \note{le 8 est à l'encre noire.}

c.\ Pour $S$ strictement hensélien et $G$ lisse, on a $H^1(S, G) = 0$. Il nous
faut donc vérifier que $H^1(X, G) = 0$. Supposons que ce soit déjà établi
après le changement de base $i : S_0 \to S$, où $S_0$ est défini par un idéal
de carré nul \add{et de longueur 1} $I$ \struck{de} de $\mathcal{O}_S$. Les
$G$-torseurs sur $X$ sont alors des déformations du $G$-torseur trivial sur
$X_0 = X \times_S S_0$. Si $L = e^{*} \Omega^1_{G/S}$ ($e$ section neutre)
ces déformations sont classifiées par
\[
\struck{H^1(S_0}
\]
\[
H^1(X_{\struck{s}}, f^{*} \underline{\operatorname{Hom}}(L, I)) = 0,
\]
et ceci prouve 12.\struck{7}9.

\page{27}
\note{Folio 12, « XVIII 1.2 ». Toute la page est barrée de deux grandes
diagonales au crayon ; elle se lit entièrement sous elles.}

Pour $S$ artinien, on a (1) $\Rightarrow$ (1') : si $G$ vérifie (1), donc est
fini localement libre, soient $G'$ son dual de Cartier, et $A(G)$ le schéma
en groupe \add{(lisse)} des éléments inversibles de l'algèbre affine de $G'$ ;
$G \simeq G''$ s'identifie au sous-groupe de $A(G)$ formé des caractères de
$G'$. Soit donc une suite exacte
\[
0 \to G \to G_1 \to G_2 \to 0
\]
avec $G_1$ et $G_2$ lisses. La suite exacte longue de cohomologie fournit, par
1.2.3 et 1.2.2 appliqué à $G_1$ et $G_2$
\[
0 \to R^1 p_{*} G \to \underline{\operatorname{Hom}}(\mathbb{G}_m, G_1) \to
\underline{\operatorname{Hom}}(\mathbb{G}_m, G_2)
\]
de sorte que $\underline{\operatorname{Hom}}(\mathbb{G}_m, G) \xrightarrow{\ \sim\ }
\struck{\operatorname{Hom}} R^1 p_{*} G$.

Dans le cas (iii), le problème est local sur $S$ (1.2.4).

Quitte à étendre $k$, on peut supposer l'existence d'une suite exacte
\[
0 \to \mathbb{G}_m^n \to G_0 \to H \to 0
\]
telle que $H$ ne contienne pas de tore. On applique alors 1.2.3 et 1.2.8 à la
suite exacte longue de cohomologie.

\section*{Structures complexes sur le disque et torseurs sous $PSL_2(\mathbb{R})$ (page 28)}

\page{28}
\note{Encre bleue sur papier quadrillé jauni, tête-bêche par rapport au
feuillet ; la page est barrée au crayon de longues diagonales et se lit sous
elles. Notes brèves, disposées en colonnes ; l'ordre suivi est celui de la
page, de haut en bas et de gauche à droite.}

$D$ : (\struck{\ill{}}

$\mathbb{R}^2$, $\omega$\quad str.\ complexe\quad $(V_0, J)$

i.e.\ \struck{$V_0$} $V = V_0 \otimes_{\mathbb{R}} \mathbb{C}$ et $\omega^{-1}$,\quad $\omega > 0$ ?

\emph{\uncertain{Gras.}}\quad $\mathbb{C}$.

$0 \to \omega \to V \to \omega^{-1} \to 0$\quad attaché à $D$

$(D, V)$ ; $V_0$, $\nabla$, $0 \to \omega \to V \to \omega^{-1} \to 0$, str.\ réelle.\note{encadré.} \emph{Aut} : $SL_2(\mathbb{R})$

« disque orienté ».

$\operatorname{Aut}(V_0)$ : $SL_2(\mathbb{R})/\pm 1$
\note{un grand trait oblique suit « $\pm 1$ », comme pour le biffer.}

\emph{en part} \quad torseur sous $PSL_2(\mathbb{R})$ (plat) sur disque abst.

$\to$ torseur plat sous $PSL_2(\mathbb{R})$ sur tt courbe\note{encadré.}

(azumaya \emph{$\operatorname{Aut}(V)$})

a/ se relève-t-il à $SL_2(\mathbb{R})$ ?

$\|$ fibré en dr.\ $\pi$ \quad $<$ (parties principales ?)\note{cerclé.}

\struck{$\mathbb{P}(P^1(X))$}
\note{la formule encadrée est barrée de traits obliques ; lecture douteuse.}

$\tilde{g}$ ok $\Rightarrow$ $k \tilde{g}$ ok ; $\Leftrightarrow$ pour $k$ impair.
\note{sous « $g$ » et « $k\,g$ » un trait vertical descend sur un $X$ :
ce sont des revêtements de $X$ ; le tilde est notre transcription de cette
disposition. Le « $g$ » de gauche est surchargé. En dessous, cerclé :
« $g-1$ ».}

Aut(rév.\ univ.\ en $x$)

obstruction\quad $H^2(X, \mathbb{Z}/2)$\quad [\emph{ou} $0$, \emph{ou} $g-1$

choix\quad $H^1(X, \mathbb{Z}/2)$

Aut\quad $H^0(X, \mathbb{Z}/2)$

\emph{$\Gamma$ de $PSL_2$ peut-il se relever à $SL_2$}

$x_1, \ldots, x_{2g}$ avec $(x_1 x_2) \cdots (x_{2g-1} x_{2g}) = e$\quad
relever $x_i$ ; \quad Rel $= \pm 1$ !
\note{les $(x_i x_{i+1})$ sont des commutateurs ; le produit est souligné
d'une accolade.}

$(x_1 x_2)(x_3 x_4) = \pm 1$

\section*{SGA 4 XVII : références à vérifier (page 29)}

\page{29}
\note{Crayon, sur papier quadrillé.}

SGA 4 XVII

\emph{Vérifier réf} : « $F$ premier aux car rés.\ de $X$ » \quad XVII 0.13

Ajouter : $X \to S$, $F$ sur $X$ premier aux car de \struck{\ill{}} $S$

: ? $\forall s \in S$, $F \mid X \times_S \operatorname{Spec}(\mathcal{O}_{S,s})$ tel que $p_s$ inv.

\emph{Ajouter} : 1.2.9 : plat / lisse par Illusie --

\section*{Une lettre signée P.~Deligne (page 30)}

\page{30}
\note{Encre bleue sur papier quadrillé ; « L 4 » en haut à droite. La
lettre, signée « P. Deligne », accompagne un envoi de pièces sur
l'exposé XVIII et le § 3.1 sur $Rf^{!}$.}

\emph{Manque} : \emph{a} Gysin \quad \emph{b} $g^{*} Rf^{!} \xrightarrow{\ \sim\ } Rf'^{!} g^{*}$ si …

\emph{c} Un grand nombre de compatibilités idiotes

\emph{d} Compatibilité aux classes fondamentales locales.

Je n'ai pas l'intention de combler ces manques dans un temps prévisible. Je
me sens toutefois une obligation morale de combler 3.1.10.3, et de donner les
compatibilités requises pour identifier $Rf^{!}$ à $f^{*}(\ )[\ ]$ (i.e.,
telle flèche définie en 3.1 est telle flèche). \uncertain{Ce} risque de faire
une vingtaine de pages illisibles et autant de jours de travail idiot pour
moi -- Peut-être le ferais-je avant que SGA 4 ne devienne un livre --
\note{ce passage, de « de combler 3.1.10.3 » à « un livre », est marqué dans
la marge gauche d'un crochet au crayon.}

Je ne me sens pas d'obligations morales pour le reste, moins emmerdant. Le
début de \emph{d} est sans doute à remonter en 1.1.

\emph{Ci-joint} : tes lettres et notes sur les jacobiennes -- une lettre à
Serre -- ta note sur Gysin --

de tes commentaires anciens et nouveaux

une photocopie d'une liste de référence dans XVII.

une liste incomplète de compatibilités manquantes.

une ancienne version, avec tes commentaires, de 3.1.

Bien à toi

P.~Deligne

\section*{Liste de compatibilités manquantes (page 31)}

\page{31}
\note{Encre bleue sur papier quadrillé, de la même main que la page 30 :
c'est vraisemblablement la « liste incomplète de compatibilités manquantes »
que la lettre annonce.}

\emph{Liste de compatibilités manquantes}

a) comp.\ de « ch base » avec $Rf_{!} \to Rf_{*}$. \struck{comp}

b) comp de « ch base » avec $f^{*}(d)[2d] \to Rf^{!}$

c) comp.\ de induction avec $f^{*}(d)[2d] \to Rf^{!}$

d) comp.\ de cl.\ fond.\ loc.\ avec $Rj^{!} Rf^{!} = R(fj)^{!}$

e) classe de correspondance

f) Th de comparaison, sur $\mathbb{C}$

g) comparaison avec la dualité des courbes par autodualité des jacobiennes

h) morphisme de Gysin et la requête de Jouanolou.

i) ch de base : transp de (\struck{\ill{}} \struck{$Rf_{!}\, Rg'_{*} \to Rg_{*}\, Rf'_{!}$}

\add{en} $R'f^{!} g^{*} \longleftarrow g'^{*} Rf^{!}$ ; cas où $f$ est lisse.
\note{la première flèche de i) est biffée d'un trait noir puis d'un trait
vert. « $R'f^{!}$ » est écrit ainsi ; on attend $Rf'^{!}$.}

\section*{« État avant Illusie » : deux fiches de calcul (pages 32 et 33)}

\page{32}
\note{Fiche à l'encre noire sur papier quadrillé, écrite dans les deux sens.
En haut, à l'encre rouge et soulignée d'un trait rouge : « ETAT Avant
Illusie ». On donne d'abord la moitié écrite à l'endroit, puis l'autre,
retournée.}

\struck{$H^1(H^0) \to H^0(H^1)$}

$X$\quad $H^1(H^0) \to H^1 \to H^0 H^1 \to H^2 H^0$

\note{sous $H^1(H^0)$ est écrit « (0) ». Plus bas, en face d'un « $S$ »,
un $H^1$ d'où partent une flèche verticale vers le $H^1$ de la ligne et une
flèche double, oblique, vers $H^0 H^1$ : c'est la suite exacte des termes
de bas degré d'une suite spectrale de Leray pour $X \to S$, écrite sans
indices ni coefficients.}

\note{Moitié retournée.}

$\langle \mathcal{L}, \mathcal{M} \rangle = \uncertain{N}_{E/S}(\mathcal{M})$\quad si $\mathcal{L} = \mathcal{O}(E)$

Tout $F(\mathcal{L})$ vient d'un $\mathcal{M}$.
\note{ces deux lignes sont réunies par une accolade.}

$\mathcal{O}(E)$ \quad $\mathcal{O}(F)$ \quad $\mathcal{O}$

$0 \to \mathcal{O}(-F) \to \mathcal{O} \to \mathcal{O}_F \to 0$

$\langle \mathcal{O}(E), \mathcal{O}(-F) \rangle = N_{E/}(\mathcal{O}(-F))\, N_{E/}(\ )^{-1}$

$\det(\ )\det(\ ) = \det(Rf_{*}(\mathcal{O}_E \otimes^{L} \mathcal{O}_F))^{-1}$
\note{sous le $\otimes$, un indice peu lisible, lu $L$.}

$\| \langle \mathcal{O}(E), \mathcal{O}(F) \rangle = \det Rf_{*}(\mathcal{O}_E \otimes \mathcal{O}_F)^{-1}$
\note{« $\|$ » est un double trait vertical dans la marge.}

$\langle 2\ 1 \rangle$ --- $\langle 1\ 1 \rangle$

$\langle \mathcal{O}_X \otimes \mathcal{O}_Y \rangle \overset{\langle f, f^{-1} \rangle}{\longrightarrow} \langle \mathcal{O}_Y \otimes \mathcal{O}_X \rangle$
\note{les deux lignes sont reliées par des traits verticaux qui descendent
chacun sur un « $0$ » ; « $\langle f, f^{-1} \rangle$ » est souligné deux
fois et tient lieu de flèche ; la lecture « $\otimes$ » dans les crochets de
gauche est incertaine (on lit plutôt « $\times$ »).}

\page{33}
\note{Fiche à l'encre noire et rouge sur papier quadrillé. Les notations
noires sont des diviseurs et un symbole modéré, les rouges des
présentations et de la descente ; on les donne dans cet ordre.}

$X_1\ Y_1$ \quad $X_1\ T$

$X_2\ Y_2$

$X_1 \cup Y_2$\quad $(-1)^{v_x f\, v_x g}\,(g^{v_x f}/f^{v_x g})$

$X_1 \cup T$\quad $(-1)^{v_x f\, v_x g}\ g^{v_x}$

$X_{\uncertain{1}} \cup Y_2$

$X_1 \cup T$\quad \struck{$T$} $\cup\, Y_2$
\note{ces deux derniers sont reliés par un trait en zigzag.}

$X_1\ Y_1$ \quad $X_1\ T$ \quad $X_3\ Y_3$ \quad $X_2\ Y_2$

$X_1\ Y_3 \mid X_3\ Y_2$

\note{Encre rouge.}

$\operatorname{Sym} T'' \quad \operatorname{Sym} T' \quad \operatorname{Sym} T$ \quad ($\operatorname{Sym} T' \to \operatorname{Sym} T$)

$T'' \rightrightarrows T' \to T$, \quad $T \to S$
\note{le premier terme est écrit $T^{1}$ ou $T''$ ; le dernier $T$ est
surchargé.}

\emph{desc.\ el.}\quad \emph{desc.\ finie.}

$A_x \to A$\quad $\equiv h.$

$M_2 \quad M_1 \quad M_0$

$A_2 \Rrightarrow A_1 \rightrightarrows A_0 \to A$
\note{avec, sous les flèches, des flèches en sens inverse (dégénérescences).}

\emph{ops} : $A_0 \leftarrow A$ \quad $\overset{?}{\Longrightarrow}$ \quad ops\quad $A_0 = A$

\section*{État primitif (page 34)}

\page{34}
\note{Encre rouge, en haut d'un feuillet quadrillé ; le reste porte au
crayon des traces effacées, illisibles.}

ETAT PRIMITIF.

\emph{contient} : ? compatibilité commut\uncertain{ation} de \uncertain{6}.
\note{le dernier signe est un 6 ou un $G$.}

\section*{Changement de base et adjonction : esquisses (page 35)}

\page{35}
\note{Encre bleue sur papier quadrillé ; toute la page est barrée de trois
longues diagonales brunes. Ce sont des schémas plus que des phrases : on
en donne les formules et on décrit les figures.}

Cat\quad $\mathcal{A}^{*}$ \quad $\mathcal{B}^{*}$ \quad $\mathcal{C}^{*}$ \quad $\mathcal{D}^{*}$
\note{quatre catégories dans deux cadres, $\mathcal{A}$, $\mathcal{B}$ en
haut, $\mathcal{C}$, $\mathcal{D}$ en bas, chacune marquée d'une
astérisque ; à gauche une flèche descendante $f^{\bullet}_{!}$, au milieu
« ok » cerclé, à droite une flèche montante suivie de « ? ». Dessous, un
carré dont les sommets sont des astérisques, la flèche verticale gauche
marquée de « ! » ; les deux lignes du bas sont indicées 0 et 1.}

$\operatorname{Hom}(f^{*}F, G)$\quad $\operatorname{Hom}(\ {}_{!}\ \ {}^{!}\ )$ \emph{ou}\quad $\operatorname{Hom}(f_{!}F, G)$\quad $\operatorname{Hom}(F, f^{!}G)$\quad $\operatorname{Hom}(f^{*}G, F)$
\note{deux carrés de flèches entre croix : dans le premier, le côté du haut
est $\operatorname{Hom}(f^{*}F, G)$ et les côtés verticaux montent, celui de
droite marqué $\operatorname{Hom}(\ {}_{!}\ \ {}^{!}\ )$ ; dans le second,
les flèches horizontales vont vers la gauche, le côté gauche est
$\operatorname{Hom}(f_{!}F, G)$ (de $F$ vers $G$), le côté droit
$\operatorname{Hom}(F, f^{!}G)$, et sous le côté du bas
$\operatorname{Hom}(f^{*}G, F)$.}

\begin{tikzcd}
\bullet \arrow[d, "f_{!}"'] & \circledcirc \arrow[l, "g^{*}"'] \arrow[d, "f_{!}"] & \bullet \arrow[r, "g_{*}"] & \bullet \\
\circledcirc & \bullet \arrow[l, "g^{*}"] & \bullet \arrow[u, "f^{!}"] \arrow[r, "g_{*}"'] & \bullet \arrow[u, "f^{!}"']
\end{tikzcd}

\note{les deux carrés sont dans un même cadre, à droite duquel on lit
« \emph{dém} : \emph{fib + cofib} » ; dans chacun une petite flèche courbe
(une 2-flèche) va d'un sommet à l'autre. Au-dessous : « $\to$ \emph{ch base}
: 2 \uncertain{dessins} ».}

\emph{adj}\quad $\operatorname{Hom}(f_{!}K, L) \sim \operatorname{Hom}(K, f^{!}L)$ \quad $/X$

$\downarrow$ \quad $\downarrow$ \quad $\sim$ \uncertain{préf}

$\operatorname{Hom}_{\vee}(f^{\vee}_{!}K, L) \sim \operatorname{Hom}(K, f^{!}_{\vee}L)$ \quad $/V$

$Rf_{*} R\operatorname{Hom}(f_{!}K, L) \sim$
\note{la ligne s'arrête là.}

\section*{3.1.6, fin : les quatre flèches de changement de base (page 36)}

\page{36}
\note{Encre bleue sur papier quadrillé, d'une écriture soignée ; les
annotations au crayon, dans les marges et entre les lignes, sont données
comme notes marginales. Ce feuillet et les suivants appartiennent à un
§ 3.1 sur le foncteur $Rf^{!}$, que continue le lot 3.}

Pour tout carré cartésien de schémas

\begin{tikzcd}
X' \arrow[r, "g'"] \arrow[d, "f'"'] & X \arrow[d, "f"] \\
S' \arrow[r, "g"] & S
\end{tikzcd}

avec $f$ compactifiable, \struck{$\mathcal{E}$} et tout faisceau
\add{d'anneaux} de torsion $\mathcal{A}$ sur $S$, l'isomorphisme de
changement de base XVII 5.2
\[
g^{*} Rf_{!} \xrightarrow{\ \sim\ } Rf'_{!}\, g'^{*}
\]
entre foncteurs de $D(X, f^{*}\mathcal{A})$ dans $D(S', g^{*}\mathcal{A})$ se
transpose en un isomorphisme
\[
(3.1.6.3)\qquad Rg'_{*}\, Rf'^{!} \xrightarrow{\ \sim\ } Rf^{!}\, Rg_{*}
\]
entre foncteurs de $D^{+}(S', g^{*}\mathcal{A})$ dans $D^{+}(X, f^{*}\mathcal{A})$.
\note{dans « 3.1.6.3 » et « 3.1.6.4 », le 6 est surchargé.}

On obtient ainsi quatre « flèches de changement de base » :
\[
(3.1.6.4)\qquad
\begin{array}{ll}
g^{*} Rf_{*} \longrightarrow Rf'_{*}\, g'^{*} & \qquad g^{*} Rf_{!} \xrightarrow{\ \sim\ } Rf'_{!}\, g'^{*} \\[1ex]
Rg^{!} Rf_{*} \xrightarrow{\ \sim\ } Rf'_{*}\, Rg'^{!} & \qquad Rg^{!} Rf_{!} \longleftarrow Rf'_{!}\, Rg'^{!}
\end{array}
\]
\note{dans les deux flèches marquées $\sim$ une pointe a été ajoutée ou
surchargée à l'autre bout, de sorte qu'elles paraissent à double sens ; on
les donne dans le sens de 3.1.6.3 et de XVII 5.2. Au crayon, une accolade
court le long de la colonne de gauche, avec une flèche qui en part et une
qui y revient.}

Les flèches de la diagonale principale sont autoadjointes, et définies pour
tout carré commutatif dès que leurs sources et buts ont un sens
\marginal{(et en tous cas quand on considère les foncteurs \uncertain{dérivés},
comme définis sur \ill{})}\marginal{des catégories $D^{+}$,} ; les
isomorphismes de la seconde diagonale sont adjoints l'un de l'autre, et
définis pour un carré \struck{commutatif} \add{cartésien} de
\struck{\ill{}}schémas, avec $f$ (resp.\ $g$) compactifiable.
\marginal{Chacun de} les flèches vérifient \struck{\ill{}} \struck{de}
\add{deux} compatibilités du type XII 4.4 pour $f$ ou $g$ composé de deux
morphismes. \marginal{Peut-on dire la raison ?}

\marginal{$g'^{*} Rf^{!} \xrightarrow{\ \sim\ } Rf'^{!} g^{*}$ \quad
\uncertain{dû} \ill{} \uncertain{sous hyp.\ de lissité}}
\marginal{($Rf^{!}$ \ill{} pour $f$ lisse \ill{}) \ill{} $g$ (loc.)
acyclique (p.\ ex.\ lisse)}
\note{ces deux dernières annotations sont au crayon, au bas du feuillet ;
la seconde est très effacée.}

\page{37}
\struck{3.1.3 Constructions du foncteur $Rf^{!}$}
\note{seule ligne du feuillet, biffée, au-dessus d'un trait.}

\section*{3.1.7 à 3.1.9 : la formule d'induction (pages 38 à 40)}

\page{38}
\marginal{Qui sont $\mathcal{A}$, $\mathcal{B}$ ?}
\textbf{3.1.7.} Pour $\underline{F}$ \struck{complexe} \add{un} \struck{des}
$f^{*}\mathcal{B}$-Module\struck{s} à droite sur $X$ et $\underline{G}$
\struck{le complexe} un $\mathcal{A},\mathcal{B}$-Bimodule à gauche
\add{sur $S$}, on dispose de
\[
(3.1.7.1)\qquad f^{\bullet}_{!}\underline{F} \otimes_{\mathcal{B}} \underline{G} \longrightarrow f^{\bullet}_{!}(\underline{F} \otimes_{\mathcal{B}} f^{*}\underline{G}).
\]
Pour $\underline{H}$ faisceau de $\mathcal{A}$-modules sur $S$, cette flèche
se transpose en
\[
(3.1.7.2)\qquad \underline{\operatorname{Hom}}_{\mathcal{A}}(f^{*}\underline{G}, f^{!}_{\bullet}\underline{H}) \longrightarrow f^{!}_{\bullet}\underline{\operatorname{Hom}}_{\mathcal{A}}(\underline{G}, \underline{H}).
\]
Cette flèche se dérive en une flèche \struck{d'induction}
\[
(3.1.7.3)\qquad R\underline{\operatorname{Hom}}_{\mathcal{A}}(f^{*}L, Rf^{!}M) \xrightarrow{\ \sim\ } Rf^{!} R\underline{\operatorname{Hom}}_{\mathcal{A}}(L, M),
\]
\marginal{($L \in \operatorname{ob}$ ?, $M \in \operatorname{ob}$ ?)}, appelée
homomorphisme d'induction, dont \marginal{on va vérifier que c'est un
isomorphisme. Je} \struck{qui} rend commutatif le diagramme suivant, pour
$K \in D^{-}(X, f^{*}\mathcal{B})$, $L \in D^{-}(S, \mathcal{A}, \mathcal{B})$
\marginal{Réf.\ ?} \struck{$\operatorname{Hom}(Rf_{!}K, R\underline{\operatorname{Hom}}_{\mathcal{A}}(f^{*}L, Rf^{!}M))$}
et $M \in D^{+}(S, \mathcal{A})$.

\begin{tikzcd}[column sep=small, row sep=large, nodes={font=\scriptsize}]
\operatorname{Hom}(K, R\underline{\operatorname{Hom}}_{\mathcal{A}}(f^{*}L, Rf^{!}M)) \arrow[r, "(1)"] \arrow[d, leftrightarrow, "\wr"'] & \operatorname{Hom}(K, Rf^{!}R\underline{\operatorname{Hom}}_{\mathcal{A}}(L, M)) \arrow[r, leftrightarrow, "\sim\ \mathrm{adj.}"] & \operatorname{Hom}(Rf_{!}K, R\underline{\operatorname{Hom}}_{\mathcal{A}}(L, M)) \arrow[d, leftrightarrow, "\wr"] \\
\operatorname{Hom}(K \otimes^{L}_{\mathcal{B}} f^{*}L, Rf^{!}M) \arrow[r, leftrightarrow, "\sim\ \mathrm{adj}"'] & \operatorname{Hom}(Rf_{!}(K \otimes^{L} f^{*}L), M) \arrow[r, "(2)"'] & \operatorname{Hom}(Rf_{!}K \otimes^{L} L, M)
\end{tikzcd}

\note{sur la page, (1) et (2) sont cerclés.}

\marginal{terminologie} Dans ce diagramme, (2) (déduit de la flèche de
\struck{changement de base} \add{projection}) est un isomorphisme (XVII 5.2).
La flèche (1) est donc un isomorphisme quel que soit $K$. On vérifie
facilement que ceci implique que (3.1.7.3) est un isomorphisme. On a obtenu

\marginal{Préciser données ($f$ comp.\ $X \to S$, $\mathcal{A}$,
$\mathcal{B}$ \ill{})}
\textbf{Proposition 3.1.8.} \emph{La flèche 3.1.7.3 est un isomorphisme : on a
la « formule d'induction »} :
\[
R\underline{\operatorname{Hom}}_{\mathcal{A}}(f^{*}L, Rf^{!}M) \xrightarrow{\ \sim\ } Rf^{!} R\underline{\operatorname{Hom}}_{\mathcal{A}}(L, M).
\]
\marginal{pour $L \in \operatorname{ob}$ \quad , $M \in \operatorname{ob}$ \quad .}

\page{39}
\textbf{3.1.9} Soient $\mathcal{A}$ et $\mathcal{B}$ \struck{d'un} deux
faisceaux d'anneaux de torsion sur $S$ et soit $u : \mathcal{B} \to
\mathcal{A}$ un homomorphisme. Prenons dans 3.1.8, pour $L$, le complexe de
$\mathcal{A},\mathcal{B}$-bimodules \struck{$\mathcal{B}$} réduit à
$\mathcal{A}$ placé en degré 0. Le foncteur
\[
\struck{E}\ \operatorname{Mod}(S, \mathcal{A}) \to \operatorname{Mod}(S, \mathcal{B}) : \underline{F} \mapsto \underline{\operatorname{Hom}}_{\mathcal{A}}(L, \underline{F}),
\]
est le foncteur de restriction des scalaires \add{$u.$} de $\mathcal{A}$ à
$\mathcal{B}$. La flèche d'adjonction, est ici un isomorphisme de foncteurs
rendant commutatif le diagramme suivant.

\begin{tikzcd}
D^{+}(X, f^{*}\mathcal{A}) \arrow[r, "u."] & D^{+}(X, f^{*}\mathcal{B}) \\
D^{+}(S, \mathcal{A}) \arrow[r, "u."] \arrow[u, "Rf^{!}"] & D^{+}(S, \mathcal{B}) \arrow[u, "Rf^{!}"']
\end{tikzcd}

\note{après chaque $D^{+}$ un mot biffé à l'encre bleue et verte,
vraisemblablement « Mod ».}

On voit que le faisceau $\mathcal{A}$ ne joue qu'un rôle bidon dans la
définition de $Rf^{!}$.
\note{dans la marge gauche, des mots au crayon presque effacés, illisibles.}

\page{40}
\struck{Prenons Soit $\mathcal{B}$}
\note{seule ligne du feuillet, biffée.}


\end{document}
